% Options for packages loaded elsewhere
\PassOptionsToPackage{unicode}{hyperref}
\PassOptionsToPackage{hyphens}{url}
\documentclass[
  ukrainian,
]{article}
\usepackage{xcolor}
\usepackage[margin=2.5cm]{geometry}
\usepackage{amsmath,amssymb}
\setcounter{secnumdepth}{-\maxdimen} % remove section numbering
\usepackage{iftex}
\ifPDFTeX
  \usepackage[T1]{fontenc}
  \usepackage[utf8]{inputenc}
  \usepackage{textcomp} % provide euro and other symbols
\else % if luatex or xetex
  \usepackage{unicode-math} % this also loads fontspec
  \defaultfontfeatures{Scale=MatchLowercase}
  \defaultfontfeatures[\rmfamily]{Ligatures=TeX,Scale=1}
\fi
\usepackage{lmodern}
\ifPDFTeX\else
\setmainfont{DejaVu Serif}
\setsansfont{DejaVu Sans}
\setmonofont{DejaVu Sans Mono}
\setmathfont{STIX Two Math}
\fi
\ifPDFTeX\else
  % xetex/luatex font selection
\fi
% Use upquote if available, for straight quotes in verbatim environments
\IfFileExists{upquote.sty}{\usepackage{upquote}}{}
\IfFileExists{microtype.sty}{% use microtype if available
  \usepackage[]{microtype}
  \UseMicrotypeSet[protrusion]{basicmath} % disable protrusion for tt fonts
}{}
\makeatletter
\@ifundefined{KOMAClassName}{% if non-KOMA class
  \IfFileExists{parskip.sty}{%
    \usepackage{parskip}
  }{% else
    \setlength{\parindent}{0pt}
    \setlength{\parskip}{6pt plus 2pt minus 1pt}}
}{% if KOMA class
  \KOMAoptions{parskip=half}}
\makeatother
\usepackage{graphicx}
\graphicspath{{/Users/kalinicenkovitalijmikolajovic/.codex/worktrees/5c15/Script_for_translate_book/Output_Final/}}
\setkeys{Gin}{width=0.82\linewidth,height=0.82\textheight,keepaspectratio}
\makeatletter
\newsavebox\pandoc@box
\newcommand*\pandocbounded[1]{% scales image to fit in text height/width
  \sbox\pandoc@box{#1}%
  \Gscale@div\@tempa{\textheight}{\dimexpr\ht\pandoc@box+\dp\pandoc@box\relax}%
  \Gscale@div\@tempb{\linewidth}{\wd\pandoc@box}%
  \ifdim\@tempb\p@<\@tempa\p@\let\@tempa\@tempb\fi% select the smaller of both
  \ifdim\@tempa\p@<\p@\scalebox{\@tempa}{\usebox\pandoc@box}%
  \else\usebox{\pandoc@box}%
  \fi%
}
% Set default figure placement to htbp
\def\fps@figure{htbp}
\makeatother
\ifLuaTeX
\usepackage[bidi=basic,shorthands=off,provide=*]{babel}
\else
\usepackage[bidi=default,shorthands=off,provide=*]{babel}
\fi
\ifLuaTeX
  \usepackage{selnolig} % disable illegal ligatures
\fi
\setlength{\emergencystretch}{3em} % prevent overfull lines
\providecommand{\tightlist}{%
  \setlength{\itemsep}{0pt}\setlength{\parskip}{0pt}}
\usepackage{bookmark}
\IfFileExists{xurl.sty}{\usepackage{xurl}}{} % add URL line breaks if available
\urlstyle{same}
\hypersetup{
  pdflang={uk},
  hidelinks,
  pdfcreator={LaTeX via pandoc}}

\author{}
\date{}

\begin{document}

\section{Определение конформного
отображения}\label{ux43eux43fux440ux435ux434ux435ux43bux435ux43dux438ux435-ux43aux43eux43dux444ux43eux440ux43cux43dux43eux433ux43e-ux43eux442ux43eux431ux440ux430ux436ux435ux43dux438ux44f}

Відображення w=\(f(z)\) комплексної площини z у комплексну площину w
називається конформним у точці \(z_{\bullet}\) площини z, якщо похідна
\(f'(z_{\bullet}) \neq 0\) . Відображення \(f(z)\) називається
конформним у регіоні D, якщо \(f'(z) \neq 0\) в кожній точці регіону D.

Наприклад, відображення \(f(z) = z^2\) конформно всюди, крім точки z =
0, оскільки \(z'(z) = 2z \neq 0\) для всіх \(z \neq 0\) . З іншого боку,
відображення \(e^z\) конформно у всій z-площині, оскільки скрізь
\(f'(z) = e^z \neq 0\) . Яке призначення конформних відображень?
Відповідь полягає в тому, що, розв'язуючи рівняння Лапласа
\(\phi_{xx} + \phi_{yy} = 0\) в певній області площини змінних x і y, ми
можемо розглядати цю площину як площину комплексної змінної z.
Розглянемо комплексне відображення w = \(f(z)\) площини z у площину w, у
якій введені координати (u, v). Відображуючи w = \(f(z\)
\phi\emph{\{xy\} + \phi}\{uy\} = 0 \()\), рівняння Лапласа
\(\varphi_{xx} + \varphi_{yy} = 0\) перетворюється на нове
диференціальне рівняння в частинних похідних, яке залежить від нових
координат і v. До --- рівняння, яке \(\phi_{xx} + \phi_{uu} = 0\)
повернеться у рівняння \(\varphi_{nn} + \varphi_{nn} = 0\) . Отже, ми
придумаємо ідею знайти конформне відображення, яке переносить область з
комплексною межею у область з простою межею (нагадаємо, що відображення
\(w=z^*\) переносить межу першого квадранта площини z у дійсну вісь
площини w).

Ось кілька прикладів, які покажуть, як можна розв'язати диференціальні
рівняння у похідних за допомогою конформних відображень.

Уравнение Лапласа в верхней полуплоскости.

Припустимо, що нам потрібно розв'язати наступну задачу Діріхле у верхній
половині видимості (рис. 47.3):

(47.1)

MATHBLK1183X

MATHBLK1184X

\begin{verbatim}
\begin{aligned} (Y & \Pi) & \varphi_{xx} + \varphi_{yy} = 0, & -\infty < x < \infty, & 0 < y < \infty, \\ (47.1) & \varphi(x, 0) = \begin{cases} 0, & |x| > 1, \\ 1, & |x| \le 1. \end{cases}
\end{verbatim}

MATHBLK1185X

\[
w = \ln\left\{\frac{z-1}{z+1}\right\}
\]

MATHBLK1186X

\[
\begin{split} S\left[u_{t}\right] &= \frac{dS\left[u\right]}{dt}, \quad S\left[u_{tt}\right] = \frac{d^{2}S\left[u\right]}{dt^{2}}, \\ S\left[u_{xx}\right] &= -\left[n\pi/L\right]^{2}S\left[u\right] + \frac{2n\pi}{L^{2}}\left[u\left(0,\ t\right) + (-1)^{n+2}u\left(L,\ t\right)\right], \\ C\left[u_{xx}\right] &= -\left[n\pi/L\right]^{2}C\left[u\right] - \frac{2}{L}\left[u_{x}\left(0,\ t\right) + (-1)^{n+1}u_{x}\left(L,\ t\right)\right]. \end{split}
\]

MATHBLK1187X M A

\pandocbounded{\includegraphics[keepaspectratio]{images/_page_357_Figure_11.jpeg}}

THBLK1188X M

\end{document}
